%%
%% This is a placeholder manuscript shell for WWW'27 (Research Track — Web Infrastructure and Agentic Systems).
%% Refresh this manuscript before final submission by syncing with docs/paper/manuscript.tex
%% from the research repository (demo-20-agent-facing-web-commoncrawl-www).
%%
\documentclass[sigconf,anonymous,review]{acmart}

\acmConference[WWW '27]{The Web Conference 2027}{May 10--14, 2027}{Dublin, Ireland}
\acmYear{2027}
\acmDOI{}
\acmISBN{}

\title{Status 200 Is Not a File: A Validated Census of the Agent-Facing Web in Common Crawl}

\author{Anonymous Author(s)}
\affiliation{%
  \institution{Anonymous Institution}
  \country{Anonymous Country}
}
\email{anonymous@example.org}

\begin{abstract}
The Web research question is whether an HTTP success stored under an agent-facing path is evidence that a site publishes an agent file. Conventions such as \texttt{/llms.txt} give language-model agents a Markdown entry point to a website, and adoption is increasingly measured from web-crawl indices, where an HTTP 200 record at the path is counted as adoption. This paper tests that inference at full scale on Common Crawl. A walk of columnar index partitions of CC-MAIN-2026-39 finds that the crawler requested \texttt{/llms.txt} on 16,029 hosts, of which the majority answered 200. A seeded sample of registered domains, stratified by Tranco tier and classified from archived bytes, shows that the 200 status is frequently an HTML soft-404: 112 of 211 cached responses are soft-404s, while only 67 are conformant agent files. The soft-404 gate is validated against live random-path probes, homepage similarity, and Common Crawl MIME detection, agreeing with independent human validation. Correcting for soft-404s substantially lowers naive adoption rates. Furthermore, crawler frontier bias understates live adoption. Usable files remain small, predominantly platform-generated, and often contain broken links. We release our open-source classification gate, stratified sample, and annotation protocol to establish rigorous standards for agent-facing web measurement.
\end{abstract}

\begin{CCSXML}
<ccs2012>
<concept>
<concept_id>10002951.10003260.10003277</concept_id>
<concept_desc>Information systems~Web mining</concept_desc>
<concept_significance>500</concept_significance>
</concept>
<concept>
<concept_id>10002951.10003260.10003261</concept_id>
<concept_desc>Information systems~Web searching and information discovery</concept_desc>
<concept_significance>400</concept_significance>
</concept>
<concept>
<concept_id>10003033.10003079.10011704</concept_id>
<concept_desc>Networks~Network measurement</concept_desc>
<concept_significance>300</concept_significance>
</concept>
</ccs2012>
\end{CCSXML}

\ccsdesc[500]{Information systems~Web mining}
\ccsdesc[400]{Information systems~Web searching and information discovery}
\ccsdesc[300]{Networks~Network measurement}

\keywords{agent-facing web, llms.txt, Common Crawl, robots.txt, prompt injection, web measurement, dual governance, autonomous agents}

\begin{document}

\maketitle

\section{Introduction}
Conventions such as \texttt{/llms.txt} have emerged to provide autonomous agents and large language models with structured, machine-readable summaries of website contents. Researchers and industry observatories frequently quantify the adoption of these agent-facing conventions by querying public web-crawl indexes, interpreting any record with HTTP response code 200 at an agent-facing URI path as verified deployment. This paper demonstrates that this foundational inference is severely flawed.

\section{Web Measurement Architecture and Crawl Corpus}
We examine the agent-facing web across two major Common Crawl monthly archives: CC-MAIN-2026-39 and the dedicated July 2026 release CC-MAIN-2026-30. Index partitions are parsed to extract candidate paths (\texttt{/llms.txt}, \texttt{/llms-full.txt}), followed by retrieval of archived WARC response records. Web domains are stratified into popularity tiers using the Tranco Top 1M list (Top 1K, Top 10K, Top 100K, Top 1M, and unranked).

\section{Soft-404 Detection and Validation Gate}
To separate genuine Markdown agent manifests from server-level fallback error pages, we construct a deterministic soft-404 classification gate. The gate inspects content-type headers, DOM structure, HTML payload cues, content-length thresholds, and semantic similarity to the host's homepage. The gate is validated against:
\begin{enumerate}
    \item Live random-path probes testing server default behavior ($\kappa > 0.80$).
    \item Home-page textual distance metrics.
    \item Common Crawl's internal Apache Tika MIME detection.
    \item A two-annotator independent human validation protocol.
\end{enumerate}

\section{Empirical Findings: Prevalence and Frontier Bias}
Among archived HTTP 200 records at \texttt{/llms.txt}, over half are HTML soft-404 error responses, particularly in high-traffic commercial strata where Single Page Application (SPA) routers return standard landing pages for non-existent endpoints. Consequently, uncorrected index queries overestimate true agent file adoption. Conversely, comparing crawl archives with crawler-independent live probes demonstrates that Common Crawl's crawl frontier under-samples emerging tail domains, introducing significant frontier bias.

\section{Payload Characteristics and Operational Hazards}
We evaluate the structural integrity of verified agent files. Median word length is small, and a high proportion of files are generated by third-party CMS plugins. Furthermore, broken hyperlinks are prevalent, and prompt injection risks arise when agent files aggregate unvalidated user-generated content.

\section{Related Work}
We situate our contribution in the literature on web measurement, robots.txt governance, Common Crawl methodology, and emerging standards for autonomous web agents.

\section{Conclusion}
HTTP status 200 cannot be taken as proof of an agent-facing file. Reliable measurement of the agent-facing web requires content-level validation and explicit awareness of crawler frontier biases.

\bibliographystyle{ACM-Reference-Format}
\bibliography{references}

\end{document}
